\uuid{1D68}
\titre{Dérivée seconde ou plus}
\auteur{bodin ; hanani ; mzari}
\organisation{exo7}
\licence{CC BY-NC-SA 4.0}
\datecreate{2020-01-20}
\niveau{L1}
\module{Analyse}
\chapitre{Dérivabilité des fonctions réelles}
\sousChapitre{Autre}
\difficulte{2}
\theme{Dérivées, dérivée seconde ou plus}
\exo7qcmid{lille-1.derivees.020}

\contenu{
\begin{question}
\qclassification{124.00}
\qtags{motcle=dérivée seconde ou plus}
Soit $f(x)=\left\{\begin{array}{ll}x^2\ln (x^2)&\mbox{si }x\neq 0\\ 0&\mbox{si }x=0.\end{array} \right.$

Quelles sont les bonnes réponses ?
\begin{answers}
    \good{$f$ est de classe $\mathcal{C}^1$ sur $\Rr$.}
    \good{$\forall x\in \Rr^*$, $\displaystyle f''(x)=2\ln x^2+6$}
    \bad{$f$ est deux fois dérivables sur $\Rr$.}
    \bad{$f$ est de classe $\mathcal{C}^2$ sur $\Rr$.}
\end{answers}
\begin{explanations}
Les théorèmes généraux assurent que $f$ est de classe $\mathcal{C}^2$ sur $\Rr^*$ avec
$$f'(x)=2x\ln (x^2)+2x\mbox{ et }f''(x)=2\ln x^2+6\mbox{ si }x\neq 0$$
et $\displaystyle \lim _{x\to 0}\frac{f(x)-f(0)}{x-0}=0=f'(0)$. On a aussi
$$\lim _{x\to 0}f'(x)=0=f'(0) \Rightarrow f'\mbox{ est continue en }0.$$
Ainsi $f$ de classe $\mathcal{C}^1$ sur $\Rr$ et est deux fois dérivable sur $\Rr^*$. Elle n'est pas deux fois dérivables en $0$ car
$$\lim _{x\to 0}\frac{f'(x)-f'(0)}{x-0}=\lim _{x\to 0}[2\ln (x^2)+2]=-\infty.$$
\end{explanations}
\end{question}
}
